%%%% CAPÍTULO 1 - INTRODUÇÃO
%%
%% Deve apresentar uma visão global da pesquisa, incluindo: breve histórico, importância e justificativa da escolha do tema,
%% delimitações do assunto, formulação de hipóteses e objetivos da pesquisa e estrutura do trabalho.

%% Título e rótulo de capítulo (rótulos não devem conter caracteres especiais, acentuados ou cedilha)

\chapter{Introdução}\label{cap:introducao}

Neste capítulo, são apresentadas as considerações iniciais do trabalho, os objetivos gerais e específicos da pesquisa, a justificativa para a escolha do tema e a estrutura do trabalho.

\section{Considerações iniciais}\label{sec:consideracoesIniciais}

A resolução de labirintos é uma tarefa clássica que demanda raciocínio espacial, planejamento e adaptação a ambientes estruturados. Tradicionalmente, algoritmos simbólicos e métodos baseados em busca foram empregados para solucionar esse tipo de problema. No entanto, com o avanço dos \textit{\gls{llms}}, surge a oportunidade de investigar como essas arquiteturas, originalmente projetadas para linguagem natural, podem ser adaptadas para tarefas que exigem compreensão espacial e manipulação de representações simbólicas.

Apesar do sucesso dos \gls{llms} em tarefas linguísticas, ainda há lacunas no entendimento de como esses modelos internalizam e representam informações espaciais, especialmente após processos de especialização como o \textit{fine-tuning}. O caráter de ``caixa-preta'' dos \gls{llms} dificulta a análise dos mecanismos internos responsáveis pelo aprendizado e pela tomada de decisão, tornando essencial o uso de técnicas de explicabilidade para investigar as transformações ocorridas nas redes neurais durante o treinamento para tarefas específicas.

\section{Objetivos}\label{sec:objetivos}

Neste contexto, este trabalho propõe-se a investigar como \gls{llms} abertos, neste estudo com 1,5 bilhões de parâmetros, aprendem a resolver labirintos representados em \gls{ascii} por meio de diferentes estratégias de \textit{fine-tuning}.

\subsection{Objetivo geral}\label{subsec:objetivoGeral}

O objetivo central é analisar as mudanças nas ativações internas das redes neurais associadas ao aprendizado da tarefa, utilizando a biblioteca \gls{llm-mri} \cite{costa2024llmmri} para visualizar e comparar as representações neuronais antes e depois do processo de aprendizagem além de avaliar métricas de redes complexas aplicadas aos grafos de ativações gerados.

\subsection{Objetivos específicos}\label{subsec:objetivosEspecificos}

Dentre os objetivos específicos deste estudo, destacam-se:

\begin{itemize}
    \item Verificar como diferentes estratégias de \textit{fine-tuning} refletem nas ativações internas dos modelos;
    \item Analisar as estruturas dos grafos de ativações recorrentes gerados a partir dos embeddings dos modelos;
    \item Avaliar métricas de grafos para identificar padrões emergentes e alterações nas representações internas; 
    \item Discutir as implicações dos resultados para o entendimento do raciocínio espacial em \gls{llms}.
\end{itemize}

\section{Justificativa}\label{sec:justificativa}

Ao abordar essas questões, este trabalho visa contribuir tanto para o avanço no entendimento prático das técnicas de especialização de \gls{llms} quanto para o aprofundamento teórico sobre os processos de representação e raciocínio em redes neurais profundas. O foco está no entendimento dos mecanismos de raciocínio espacial em \gls{llms}, avaliando desempenho e interpretabilidade dos modelos em tarefas estruturadas de navegabilidade em labirintos textuais.

Como parte dos resultados obtidos, observou-se que as diferenças nas ativações internas dos modelos após o \textit{fine-tuning} foram menos discrepantes do que inicialmente esperado. Embora o \textit{fine-tuning} influencie o processamento de informações nos modelos, as mudanças estruturais nos grafos de ativações não foram significativamente refletidas nas métricas analisadas. Essas observações são discutidas em detalhes na seção de discussão, onde são exploradas possíveis razões para esses achados e suas implicações para futuras pesquisas na área.

\section{Estrutura do Trabalho}\label{sec:estruturaTrabalho}

O restante desta proposta está organizado da seguinte forma: o \autoref{cap:referencialTeorico} detalha os fundamentos teóricos necessários para o entendimento do tema; o \autoref{cap:trabalhos:relacionados} evidencia os trabalhos relacionados relevantes; o \autoref{cap:metodologia} detalha a metodologia que foi adotada para o desenvolvimento deste estudo, incluindo a preparação dos dados, seleção de modelos, estratégias de ajuste e análise das métricas sobre os grafos de ativações neurais; os resultados obtidos são discutidos no \autoref{cap:resultados}; por fim, o \autoref{cap:conclusoeseperspectivas} apresenta as conclusões, limitações e perspectivas para o aprofundamento do estudo.

